
Model repositories contain millions of neural network checkpoints, but metadata is often corrupted or missing. This work investigates whether neural network weights encode task identity across diverse architectures, enabling inference without metadata. We introduce a Hierarchical Variational Autoencoder that embeds heterogeneous architectures into unified weight space by trading local equivariance for global expressivity. The three-level design combines architecture-specific equivariant encoders, permutation-invariant pooling, and Transformer sequence modeling. Proof-of-concept validation on synthetic model zoo data containing 2,120 models spanning 2 architecture families (CNNs with depth variants: small, large; ResNets with variants: 18, 34) and 9 vision tasks (CIFAR-10, CIFAR-100, TinyImageNet, MNIST, FashionMNIST, SVHN, USPS, STL10, EuroSAT) demonstrates that same-task different-architecture models cluster significantly tighter than random baselines (Within-Cluster Sum of Squares ratio 0.495, p<0.000001, Cohen's d=1.45). Architecture subspaces exhibit strong compatibility (CKA similarity 0.82 for same-task pairs vs 0.14 for different-task pairs), and hierarchical pooling preserves 68\% task-relevant information. These synthetic data results suggest that task-level functional constraints may create architecture-invariant structural features in weight distributions. Real dataset validation on ModelZooDataset checkpoints is required before publication to confirm external validity and rule out synthetic data artifacts.

